\chapter{Introduction}
\label{section:introduction}

Modern computing workloads are rarely executed on a single machine. A useful system must accept many tasks, understand what resources each task needs, select a suitable worker, run the task safely, and store the result. This becomes harder when workers are heterogeneous. One worker may have more CPU, another may have more memory, and the best choice can change every few seconds as tasks start and finish.

The project developed in this work is called \textbf{Agentic Cloud Cluster}. It is a Go-based distributed cluster managing framework that schedules Docker-based tasks across worker nodes. The focus of the project is the cluster framework and the PPO-based scheduler. Other surrounding utilities are treated only as supporting parts of the system.

\section{Project Definition}
\label{section:proj_def}

The project can be defined as follows:

\begin{quote}
To build a distributed task execution framework in Go and improve task placement decisions using a reinforcement learning scheduler based on Proximal Policy Optimization.
\end{quote}

The system has a master node and multiple worker nodes. The master accepts tasks, tracks workers, stores task state, and asks a scheduler to select a worker. The worker executes the task using Docker containers and reports status back to the master. The scheduling algorithm can be changed without changing the worker execution logic.

\section{Motivation}

The first scheduling solution most systems use is a simple rule. Round-Robin is easy to understand: send one task to worker 1, the next to worker 2, and so on. This is useful as a baseline, but it ignores whether the worker has enough resources, whether it is already busy, and whether the task is likely to miss its deadline.

Heuristic schedulers improve this by adding rules. For example, a scheduler can prefer the worker with the most available CPU or the lowest estimated risk. However, fixed rules are difficult to tune because the right trade-off changes with workload type. A CPU-heavy task, a memory-heavy task, and a burst of short tasks do not need the same decision rule.

This is why reinforcement learning was explored. In reinforcement learning, an agent learns by taking actions and receiving rewards \cite{sutton2018reinforcement}. For cluster scheduling, the state is the current cluster condition, the action is the selected worker, and the reward measures whether the placement was useful. This matches the nature of scheduling: every placement changes the next state of the cluster.

\section{Project Objectives}
\label{section:proj_obj}

The objectives of the project were:

\begin{itemize}
    \item Build a working Go master-worker cluster that can run Docker-based tasks.
    \item Maintain worker registration, heartbeat, resource tracking, task lifecycle, and result storage.
    \item Provide a clean scheduling boundary so different scheduling algorithms can be compared.
    \item Implement baseline scheduling and PPO-based scheduling.
    \item Train PPO using realistic cluster traces and a reward function aligned with feasibility, queue pressure, tail pressure, and load balance.
    \item Evaluate the system with repeatable campaigns and explain the results honestly.
\end{itemize}

\section{Scope}

The core scope is the distributed cluster and PPO scheduler. The report intentionally does not focus on auxiliary interfaces or monitoring tools. They may exist as supporting utilities in the repository, but the technical contribution is the Go control plane, worker execution layer, and reinforcement learning scheduler.

\section{Main Contributions}

The main contributions are:

\begin{enumerate}
    \item A Go-based distributed task execution framework with master and worker nodes.
    \item A scheduler interface that supports simple baselines and PPO scheduling.
    \item A PPO service that receives scheduling requests over gRPC and returns a worker decision.
    \item A trace replay training environment with lifecycle-based resource accounting.
    \item A reward design that penalizes infeasible placement, queue pressure, tail pressure, and imbalance.
    \item A benchmark setup for comparing scheduling behavior across load patterns.
\end{enumerate}
